% HISTORICAL assembly fragment; the current integrated and checked text is
% parts/bmodels/bmodels.tex, which the master actually includes.
\begin{theorem}[Proved: faithfulness and the full powerset support lattice]\label{SA21-BM-cell-support}
The cell atelier is faithful and spatial. Its complete
support ortholattice for primitive internal disjointness is
\[
 \Sigma_M\cong\mathcal P(X).
\]
The open support of \(m\) corresponds to \(m^\circ\),
and its whole support to \(|m|\). This is the full
powerset, not only the definable subsets.
\end{theorem}

\begin{proof}
For arbitrary \(Q\subseteq X\), put
\(S_Q=\{m\in M:m^\circ\subseteq Q\}\).
For any \(\mathcal U\subseteq M\),
\[
 \mathcal U^\perp
 =S_{X\setminus\bigcup_{m\in\mathcal U}m^\circ}.
\]
Every point of \(Q\) supplies its singleton
multistrate in \(S_Q\), so
\[
 \bigcup_{m\in S_Q}m^\circ=Q,\qquad
 S_Q^\perp=S_{X\setminus Q},\qquad
 S_Q^{\perp\perp}=S_Q.
\]
Every closed polar has this form.
Thus \(Q\mapsto S_Q\) is a bijection
preserving inclusion and complement,
hence all complete Boolean operations.
This includes \(X=\varnothing\).
Nondefinable \(Q\) are permitted because
polarity allows arbitrary multistrate
families.

Distinct lieux are disjoint, and
primitive disjointness is exactly
disjointness of their open shadows,
proving faithfulness.
For spatiality, the ordered deployment
\(\{S_{n^\circ}:n\leq m\}\)
recovers each cell \(n^\circ\), since
\(Q\mapsto S_Q\) is injective.
Downward unions recover each closed
support, hence the whole elementary
figure \(m\). This proves
At spat 3 of 156-9, scan 50.
The incidence order and symmetric
irreflexive relation give At spat 1.
Distinct incident strata have
disjoint cells, giving At spat 2.
\end{proof}

\begin{theorem}[Proved: coarsest and elementary extensions]\label{SA21-BM-cell-extension}
If \(G\leq F\), \(G'\preceq G\), the closed-stratum family
\[
 E=G'\cup(F\setminus G)
\]
is a figure with \(E\preceq F\), \(E|G=G'\).
Every other extension \(H\preceq F\) inducing \(G'\)
refines \(E\), so \(E\) is the unique coarsest extension.
Its carrier is an order isomorphism outside the
prescribed boundary and has the required mixed
incidences. Every extension with this complement
property is equal to \(E\).
For elementary \(m\), a subdivision of \(\partial m\)
extends uniquely to an elementary \(m'\preceq m\)
with that boundary. Thus At 14, geometric
At 14 bis and At 14 ter hold.
\end{theorem}

\begin{proof}
No \(A\in F\setminus G\) lies in \(|G|\),
since its open cell is disjoint from \(|G|\).
Hence it lies below no stratum of \(G'\).
The residual cell of every new stratum
is therefore unchanged. For an old
\(A\in F\setminus G\), the old boundary
portion in \(G\) is \(A\cap|G|\).
Its induced subdivision by \(G'\)
covers exactly the same set.
Other old boundary strata remain,
so the residual cell of \(A\)
also remains unchanged.
The resulting cells are the
cells of \(G'\) and the old
cells outside \(G\).
They partition \(|F|\), retain
their stated closures and obey
the frontier condition by this
boundary calculation. They
refine \(F\), proving the claims
about \(E\).

For another extension \(H\),
an \(H\)-cell with carrier
in \(G\) belongs to \(H|G=G'\);
every other cell lies in an
old unchanged cell outside
\(G\). Thus \(H\ll E\).
Antisymmetry proves uniqueness
of a coarsest extension.
On the outside strata the
carrier of \(E\) is the
identity. For \(B\in G'\)
and \(A\in F\setminus G\),
\[
 B\subseteq A\quad\Longleftrightarrow\quad c(B)\subseteq A.
\]
One implication uses \(B\subseteq c(B)\).
For the other, the old carrier
cell of \(B\) meets \(A\);
frontier compatibility puts
its closure in \(A\).
These are the mixed incidences.
If another extension has a
bijective outside carrier,
only one cell covers each
old cell outside \(G\).
Equal support forces that
cell to be the whole old
cell, with its old closure.
Its inside figure is specified,
so the extension is \(E\).

For elementary \(m\), use
\(G=\partial m\).
The construction retains
the one top closed stratum
\(|m|\), hence is elementary
with top cell
\(|m|\setminus|\partial m|=m^\circ\).
Any other elementary extension
with the given boundary has
the same support and boundary
support, hence the same top
cell and all the same strata.
This proves uniqueness.
No triangulation is needed;
allowing general cell closures
rather than simplices is
what makes the construction
possible.
\end{proof}

\begin{counterexample}[Refuted: literal all-target At 14 bis]\label{SA21-BM-14bis-reading}
The literal mixed-order clause of 156-6, scan 88,
as transcribed and visually checked in Part B.1,
is false in the cell model if its second variable
is allowed to range over all of \(\widetilde F'\).
The corrected clause restricts that variable
to \(\widetilde F'\setminus\widetilde G'\).
\end{counterexample}

\begin{proof}
Take \(G=F\) the interval figure on \([0,1]\),
with the two endpoint cells and the open interval.
Subdivide at \(1/2\) to get \(G'\).
An extension inducing \(G'\) must be \(F'=G'\),
since \(G=F\). The two distinct open half-interval
cells have the same old interval carrier.
Literal order reflection for all pairs
of boundary strata would make each
of their closed strata lie below the other,
which is false. The five-stratum carrier
poset is completely specified and checked
by the program. The corrected mixed
condition, with the second stratum
outside \(G'\), is proved by the
previous theorem, as are the other
cartesian and complement conditions.
Thus the literal reading and its
geometric correction have separate
truth values.
\end{proof}

\begin{theorem}[Known input and proved consequence: moderation]\label{SA21-BM-cell-moderation}
Every finite collection of subdivisions of a cell figure
has a common subdivision. Every closed definable
\(S\subseteq X\) has a cell figure, and any finite
collection of its finite cell stratifications
has a common refinement. In particular the
atelier is modérée in the sense of 156-4, scan 70.
\end{theorem}

\begin{proof}
The input is compatible definable triangulation:
if \(A\subseteq R^n\) is closed, bounded and
definable and \(B_1,\ldots,B_r\subseteq A\)
are definable, there are a finite simplicial
complex \(K\) and a definable homeomorphism
\(h:|K|\to A\) such that every \(B_i\)
is a union of images of open simplices.
We use the fully opened statement of
\citet[Theorem 4.4, \S4.3, p.~42;
proof pp.~42--44]{Coste1999}, over an
arbitrary fixed o-minimal expansion
of a real closed field.
The standard book locator is
\citet[Chapter 8, \S2, Theorem 2.9]{vandenDries1998};
its theorem page was not obtained
in full and is not the evidential
basis of this proof.

Use the definable homeomorphism
\[
 b:R^n\longrightarrow(-1,1)^n,\qquad
 b(t)_i=\frac{t_i}{1+|t_i|},
\]
whose inverse is \(s_i/(1-|s_i|)\).
Let \(A=\overline{b(X)}^{\,R^n}\).
Triangulate it compatibly with
\(b(X)\), \(b(S)\), and the images
of all cells in the given finite
collection of stratifications.
For every open simplex \(\sigma\)
whose image lies in \(b(S)\), set
\[
 C_\sigma=b^{-1}(h(\sigma)).
\]
It is definably homeomorphic to
an open simplex, hence to
\((0,1)^{\dim\sigma}\) by a
semialgebraic homeomorphism.
The selected cells partition \(S\).

The closure of \(h(\sigma)\)
in \(b(X)\) is
\(h(\overline\sigma)\cap b(X)\).
An open face meeting \(b(X)\)
lies wholly in it by compatibility.
Since \(b(S)\) is closed in
\(b(X)\), every such face lies
in \(b(S)\). Thus each cell
closure in \(X\) is a union
of selected full cells, with
the frontier condition.
The closed strata form a
figure \(H\) in our exact
model. Compatibility with
all old cells makes every
new cell lie in one old
cell. It follows that
\(H\) refines all the
given figures and has
their support \(S\).
This handles unbounded
or nonclosed \(X\), and
\(S\) closed only in \(X\).
If \(S=\varnothing\), use
the empty figure.

No common linear subdivision
of fixed triangulating maps
is needed. Compatibility
with finitely many definable
cells is sufficient; no
Hauptvermutung or uniqueness
theorem is used.
\end{proof}

\begin{proposition}[Proved: the correct connectedness statement]\label{SA21-BM-cell-connected}
The undirected refinement components of \((M,\ll)\)
are exactly the definably connected components
of \(X\): a multistrate is assigned the unique
component containing its support, and two
multistrates are in the same refinement
component exactly when those assignments
agree. Different refinement components
are compatible. For an individual finite
figure, its incidence components are
exactly the definably connected components
of its support.
\end{proposition}

\begin{proof}
Generalized cells are definably connected,
since open boxes are definably path
connected. A closure of a definably
connected set is definably connected:
a definable clopen partition would
restrict to the dense set, and
density excludes an unused nonempty
open part. Thus \(|m|\) lies in
one component. Refinement implies
support inclusion, so the assigned
component is constant on refinement
edges. Each \(p\in|m|\) also
gives \(\{p\}\ll m\).

Conversely
\citet[Theorem 3.9, p.~28]{Coste1999}
gives finitely many definably
connected components, each definably
arcwise connected. Join \(p,q\)
in one component by a continuous
definable \(\gamma:[0,1]\to X\);
no injectivity is needed.
Its image \(P\) is bounded by
the extreme-value property
applied to its coordinates,
and closed by the same property
applied to distance from a
proposed missing closure point.
These follow from
\citet[Theorem 3.4, pp.~26--27]{Coste1999}.
The image is definably connected
as a continuous image of \([0,1]\).

Triangulate \(P\) compatibly
with its endpoints. The
resulting finite cell figure
in \(X\) has \(p,q\) as
vertices and connected
incidence graph: otherwise
an incidence component
and its complement would
be finite unions of closed
strata, giving a nontrivial
definable clopen partition.
Incidence implies refinement,
so this graph connects the
points and then the original
multistrates. Distinct
components have disjoint
supports, hence compatibility.

For an individual figure,
incidence components are
definably clopen by the
same finite-union argument.
They are connected because
each closed stratum is
connected and incident
strata intersect. Their
connected finite incidence
graph gives connected union,
successively along a spanning
tree. These are therefore
exactly its support components.
\end{proof}

\begin{proposition}[Refuted stronger assertions and proved scope limits]\label{SA21-BM-cell-limits}
Global incidence C4 and C4 bis, the polyhedral
restrictions and an infinite-family strengthening
of moderation fail in general. Allowing arbitrary
nonclosed supports with closures taken only
relative to themselves also destroys the
disjoint-support compatibility axiom.
\end{proposition}

\begin{proof}
For C4 let \(X=R\). The one-cell
figure on \(R\) has no proper
incidence substratum. It cannot
occur properly below another
elementary figure since its
support is already all of \(R\).
It is isolated in the global
incidence poset. The singleton
at \(0\) lies in a different
component and is incompatible
with it, since the open cells
intersect. C4 fails, and so
does its strengthening C4 bis.

For the polyhedral restrictions,
partition the semialgebraic unit
circle into two points \(p,q\)
and two complementary open arcs.
The arc closures intersect in
the two-point figure, neither
empty nor elementary. This
refutes At pol 1 and the
biconditional At \(9''\).
It does not refute CL(pol)
or At CD2\( '\), which
are sufficient criteria.
Partitioning
the circle instead into
one point and its complement
gives a top stratum with
exactly one immediate predecessor,
refuting At simp 1.
The cells are generalized
one-cells: stereographic
projection identifies a
punctured circle with \(R\),
then with \((0,1)\); the
two arcs are semialgebraically
homeomorphic to intervals.
The finite incidence tests
are in the program; these
geometric realizations put
the checked posets in the
specified model.

For arbitrary nonclosed supports,
take the one-cell partition
of \((0,1)\subset R\) and the
singleton at \(1\).
Their supports are disjoint,
but they cannot both be
closed subfigures of a
common frontier stratification:
\((0,1)\) is not closed in
\((0,1]\), so its cell
acquires the new boundary.
This is the obstruction
removed by fixed ambient
closures and closed supports.

Finally subdivide the standard
interval figure at every
\(t\in(0,1)\), one subdivision
for each \(t\). A common
refinement with the same
support would have every
singleton \(t\) as a cell,
since it must cover each
point carrier. It would
be infinite, contrary to
the figure definition.
Moderation only quantifies
over pairs. Disconnected
supports cause no problem:
their finitely many definable
components can be stratified
separately and united.
Disconnected sets are
not called individual
generalized cells.
\end{proof}
